%!TEX root = TheoreticalPhysicsC.tex
\section{Lecture: Angular Momentum}

\subsection{Orbital Angular Momentum}
After the three solvable models comes the hydrogen atom: a particle in a central potential.

We now study quantum systems which display an orbital angular momentum in the classical limit. The canonical example is that of a particle able to move in three spatial dimensions. We shall therefore assign the Hilbert space $\mathcal{H} \equiv L^2(\mathbb{R}^3)$ to model the states. Our aim is to work out the right quantum observables and their properties.

Following our quantisation table, we identify the operators below as the observables belonging to the classical angular momentum components.
\begin{definition}
	The \emph{orbital angular momentum operators} are $\widehat{\mathbf{L}} \equiv \widehat{\mathbf{r}}\times \widehat{\mathbf{p}}$, where 
	\begin{equation}
		\mathbf{r} \equiv \begin{pmatrix}
			\widehat{x} \\
			\widehat{y} \\
			\widehat{z} 
		\end{pmatrix}\quad
		\mathbf{p} \equiv \begin{pmatrix}
			\widehat{p}_x \\
			\widehat{p}_y \\
			\widehat{p}_z 
		\end{pmatrix}\quad\text{and}\quad \widehat{\mathbf{L}} \equiv \begin{pmatrix}
			\widehat{L}_x \\
			\widehat{L}_y \\
			\widehat{L}_z 
		\end{pmatrix} \equiv \begin{pmatrix}
			\widehat{L}_1 \\
			\widehat{L}_2 \\
			\widehat{L}_3 
		\end{pmatrix}.
	\end{equation}
	The position-basis matrix elements read  
	\begin{align}
		\langle x,y,z|\widehat{L}_x|\Psi\rangle &= -i{\hbar} \left(y\frac{\partial}{\partial z} - z\frac{\partial}{\partial y} \right)\Psi(x,y,z) \\
		\langle x,y,z|\widehat{L}_y|\Psi\rangle &=  i{\hbar} \left(x\frac{\partial}{\partial z} - z\frac{\partial}{\partial x} \right)\Psi(x,y,z) \\
		\langle x,y,z|\widehat{L}_z|\Psi\rangle &=  -i{\hbar} \left(x\frac{\partial}{\partial y} - y\frac{\partial}{\partial x} \right)\Psi(x,y,z),
	\end{align}
	where $\Psi(x,y,z) = \langle x,y,z|\Psi\rangle$ is an arbitrary wave function. 
\end{definition}

One can show the following commutators:
\begin{equation}
	\eqbox{[\widehat{L}_1,\widehat{L}_2] = i\hbar \widehat{L}_3, \quad [\widehat{L}_2,\widehat{L}_3] = i\hbar \widehat{L}_1\quad\text{and}\quad [\widehat{L}_3,\widehat{L}_1] = i\hbar \widehat{L}_2.}
\end{equation}
It also follows that
\begin{equation}
	\eqbox{[\widehat{\mathbf{L}}^2, \widehat{L}_j]=0, \quad \text{for all $j=1,2,3$,}}
\end{equation}
where
\begin{equation}
	\widehat{\mathbf{L}}^2 \equiv \widehat{L}_1^2+\widehat{L}_2^2+\widehat{L}_3^2.
\end{equation}

The whole commutator hangs on a single point: $\frac{\partial}{\partial z}$ commutes with $x$, but not with $z$; forget the product rule and you get zero.

The commutators are proved as follows. Let $\Psi(x,y,z)$ be an arbitrary wave function and consider
\begin{equation}
	\begin{split}
		\langle x,y,z|\widehat{L}_1\widehat{L}_2|\Psi\rangle &= {\hbar}^2 \left(y\frac{\partial}{\partial z} - z\frac{\partial}{\partial y} \right)\left(x\frac{\partial}{\partial z} - z\frac{\partial}{\partial x} \right)\Psi(x,y,z) \\
		&= {\hbar}^2 \left(xy\frac{\partial^2}{\partial z^2} - y\frac{\partial}{\partial x} - yz\frac{\partial^2}{\partial x \partial z} - xz\frac{\partial^2}{\partial y\partial z}  + z^2\frac{\partial^2}{\partial x\partial y} \right)\Psi(x,y,z)
	\end{split}
\end{equation}
and 
\begin{equation}
	\begin{split}
		\langle x,y,z|\widehat{L}_2\widehat{L}_1|\Psi\rangle &= {\hbar}^2 \left(x\frac{\partial}{\partial z} - z\frac{\partial}{\partial x} \right)\left(y\frac{\partial}{\partial z} - z\frac{\partial}{\partial y} \right)\Psi(x,y,z) \\
		&= {\hbar}^2 \left(xy\frac{\partial^2}{\partial z^2} - x\frac{\partial}{\partial y} - xz\frac{\partial^2}{\partial y \partial z} - yz\frac{\partial^2}{\partial x\partial z}  + z^2\frac{\partial^2}{\partial x\partial y} \right)\Psi(x,y,z).
	\end{split}
\end{equation}
We therefore compute the commutator
\begin{equation}
	\langle x,y,z|\widehat{L}_1\widehat{L}_2 - \widehat{L}_2\widehat{L}_1|\Psi\rangle = {\hbar}^2 \left(- y\frac{\partial}{\partial x} + x\frac{\partial}{\partial y} \right)\Psi(x,y,z) = i\hbar \langle x,y,z|\widehat{L}_3|\Psi\rangle.
\end{equation}
Since this holds for every $|\Psi\rangle$, it is an operator identity.

That $\widehat{\mathbf{L}}^2$ and $\widehat{L}_3$ possess common eigenstates is a special case of a general fact.
\begin{lemma}
	If $\widehat{A},\widehat{B}:\mathcal{H}\to\mathcal{H}$ are Hermitian with $[\widehat{A},\widehat{B}]=0$, then they possess common eigenvectors.
\end{lemma}
\begin{proof}
	Finite dimension, no degeneracy: from $\widehat{A}|\phi_j\rangle = a_j|\phi_j\rangle$ and $|\psi_j\rangle \equiv \widehat{B}|\phi_j\rangle$ it follows that
	\begin{equation}
		\widehat{A}|\psi_j\rangle = \widehat{A}\widehat{B}|\phi_j\rangle = \widehat{B}\widehat{A}|\phi_j\rangle = a_j\widehat{B}|\phi_j\rangle = a_j|\psi_j\rangle.
	\end{equation}
	The eigenspace belonging to $a_j$ is one-dimensional, so $|\psi_j\rangle \propto |\phi_j\rangle$, i.e.\ $\widehat{B}|\phi_j\rangle = b_j|\phi_j\rangle$. If there is degeneracy one diagonalises $\widehat{B}$ inside the eigenspace; the statement holds in general too, but is harder to justify.
\end{proof}
The common eigenstates of $\widehat{\mathbf{L}}^2$ and $\widehat{L}_3$ are denoted $|\ell,m\rangle$ and satisfy the following eigenvalue equations 
\begin{equation}
	\eqbox{
	\begin{aligned}
		\widehat{\mathbf{L}}^2|\ell,m\rangle &= \hbar^2\ell(\ell+1)|\ell,m\rangle, \\
		\widehat{L}_3|\ell,m\rangle &= \hbar m|\ell,m\rangle,
	\end{aligned}
	}
\end{equation}
where $\ell = 0, \frac12, 1, \frac32, \ldots$ and $m = -\ell, -\ell + 1, \ldots, \ell$. Here $\ell$ is called the \emph{angular momentum quantum number} and $m$ the \emph{magnetic quantum number}.

The eigenvalues are obtained just as for the harmonic oscillator. With
\begin{equation}
	\eqbox{\widehat{L}_\pm \equiv \widehat{L}_1 \pm i\widehat{L}_2}
\end{equation}
the commutation relations above give
\begin{equation}
	[\widehat{L}_3,\widehat{L}_\pm] = \pm\hbar\widehat{L}_\pm, \quad [\widehat{L}_+,\widehat{L}_-] = 2\hbar\widehat{L}_3 \quad\text{and}\quad [\widehat{\mathbf{L}}^2,\widehat{L}_\pm] = 0.
\end{equation}
The $\widehat{L}_\pm$ act as ladder operators: they shift the eigenvalue of $\widehat{L}_3$ by $\pm\hbar$ and leave that of $\widehat{\mathbf{L}}^2$ untouched. Since $\widehat{L}_3^2$ can never grow larger than $\widehat{\mathbf{L}}^2$, the ladder breaks off at the top and at the bottom, and that forces the list above.

% fig_vl14_multiplett.tex — Das Drehimpuls-Multiplett: die Leiter in m und
% der Kegel fester Länge, beide für ell = 2.
%
%   * geometry comes from the eigenvalue equations derived in the lecture,
%     never from a freehand sketch:
%         L^2 |ell,m> = hbar^2 ell(ell+1) |ell,m>,
%         L_3 |ell,m> = hbar m |ell,m>,   m = -ell, ..., +ell,
%         [L_3, L_pm] = +- hbar L_pm,     L_+|ell,ell> = L_-|ell,-ell> = 0.
%     Both panels are drawn in units of hbar, with ell = 2, so
%     |L|/hbar = sqrt(ell(ell+1)) = sqrt6 = 2.4495 is a fixed length and
%     the rungs sit at the integers -2,...,2.  The rung spacing is 1
%     because [L_3,L_pm] = +- hbar L_pm; nothing is spaced by eye.
%
%   * the point of the plate is the number sqrt6 = 2.4495 against the
%     number 2.  The top rung of the ladder falls SHORT of the length of
%     the vector, and that gap is what forces the ladder to break off and
%     what keeps the vector off the axis.  Panel (a) shows the gap as a
%     distance, panel (b) shows the same gap as an angle:
%         cos theta_min = ell / sqrt(ell(ell+1)) = sqrt(ell/(ell+1)),
%     i.e. 35 deg at ell = 2 and only 6 deg at ell = 100 — the classical
%     limit, as an order of magnitude rather than as an assertion.
%
%   * in panel (b) the vector of state m is drawn to the right rim of its
%     cone, at (sqrt(ell(ell+1)-m^2), m).  That point lies exactly on the
%     circle of radius sqrt6: the arrows END on the dashed circle by
%     construction, they are not fitted to it.  The cones are ellipses
%     with vertical semi-axis 0.15 times the horizontal one — the flattest
%     perspective that still reads as a cone, chosen because at 0.24 the
%     m = 0 and m = 1 cones would already overlap near the axis.
%
%   * colour encodes the magnetic quantum number m and nothing else: one
%     m, one colour, both panels, so the top rung in (a) and the top cone
%     in (b) are the same state.  The ramp runs m = -2,...,+2 through
%     `tpclevels` (spacecadet, cgblue, munsell, birrenorange, birrenred).
%     Ladder arrows and construction arcs are neutral: the operators are
%     not quantum numbers.
%
%   * the region |L_3| > |L|, which no state can reach, is hatched rather
%     than filled, so it survives greyscale printing; the bound itself is
%     a neutral grey rule, a reference line rather than data.  Following
%     fig_vl11_kasten, the hatch is a band beside that rule rather than a
%     full-bleed field, so the two L_pm|ell,+-ell> = 0 labels stand on
%     clean paper instead of on the hatch.
%
% Panel (a) the 2ell+1 rungs with L_+ and L_- between them, and the two
%           blocked arrows at the ends where the ladder terminates.
% Panel (b) the same five states as cones of fixed slant length: the
%           vector never lies along L_3, and the minimum tilt shrinks
%           only slowly with ell.

\begin{center}
\begin{tikzpicture}

% ─────────────────────────────────────────────────────────────
%  (a)  The ladder in m:  L_3 |ell,m> = hbar m |ell,m>
% ─────────────────────────────────────────────────────────────
\begin{axis}[
    name=ladder,
    tpclevelplot,
    scale only axis,
    width=5.4cm, height=5.8cm,
    xmin=-1.75, xmax=1.45,
    ymin=-3.6,  ymax=3.6,
    title={(a) $\ell=2$: \tpcifenglish{multiplet and ladder}%
                                      {Multiplett und Leiter}},
]

  % |L_3| > |L| is unreachable for every state: hatched, not filled.
  \fill[forbidden] (axis cs:-1.75,2.4495)  rectangle (axis cs:1.45,3.16);
  \fill[forbidden] (axis cs:-1.75,-2.4495) rectangle (axis cs:1.45,-3.16);

  % The bound itself, a neutral reference rule at |L|/hbar = sqrt6.
  \draw[levellinegrey] (axis cs:-1.75,2.4495)  -- (axis cs:1.45,2.4495);
  \draw[levellinegrey] (axis cs:-1.75,-2.4495) -- (axis cs:1.45,-2.4495);

  % The 2ell+1 = 5 rungs, at L_3/hbar = m.  Colour = m.
  \draw[line width=1.1pt, birrenred]    (axis cs:-0.75,2)  -- (axis cs:0.65,2);
  \draw[line width=1.1pt, birrenorange] (axis cs:-0.75,1)  -- (axis cs:0.65,1);
  \draw[line width=1.1pt, munsell]      (axis cs:-0.75,0)  -- (axis cs:0.65,0);
  \draw[line width=1.1pt, cgblue]       (axis cs:-0.75,-1) -- (axis cs:0.65,-1);
  \draw[line width=1.1pt, spacecadet]   (axis cs:-0.75,-2) -- (axis cs:0.65,-2);

  \node[lbl, anchor=east] at (axis cs:-0.82,2)  {$m=+2$};
  \node[lbl, anchor=east] at (axis cs:-0.82,1)  {$m=+1$};
  \node[lbl, anchor=east] at (axis cs:-0.82,0)  {$m=0$};
  \node[lbl, anchor=east] at (axis cs:-0.82,-1) {$m=-1$};
  \node[lbl, anchor=east] at (axis cs:-0.82,-2) {$m=-2$};

  % L_+ raises m by one:  [L_3,L_+] = +hbar L_+
  \draw[-{Stealth[length=5pt,width=3.5pt]}, spacecadet, line width=0.9pt]
    (axis cs:0.42,-1.86) -- (axis cs:0.42,-1.14);
  \draw[-{Stealth[length=5pt,width=3.5pt]}, spacecadet, line width=0.9pt]
    (axis cs:0.42,-0.86) -- (axis cs:0.42,-0.14);
  \draw[-{Stealth[length=5pt,width=3.5pt]}, spacecadet, line width=0.9pt]
    (axis cs:0.42,0.14)  -- (axis cs:0.42,0.86);
  \draw[-{Stealth[length=5pt,width=3.5pt]}, spacecadet, line width=0.9pt]
    (axis cs:0.42,1.14)  -- (axis cs:0.42,1.86);
  \node[lbl, anchor=west] at (axis cs:0.50,1.5) {$\widehat{L}_+$};

  % L_- lowers m by one:  [L_3,L_-] = -hbar L_-
  \draw[-{Stealth[length=5pt,width=3.5pt]}, spacecadet, line width=0.9pt]
    (axis cs:-0.52,1.86)  -- (axis cs:-0.52,1.14);
  \draw[-{Stealth[length=5pt,width=3.5pt]}, spacecadet, line width=0.9pt]
    (axis cs:-0.52,0.86)  -- (axis cs:-0.52,0.14);
  \draw[-{Stealth[length=5pt,width=3.5pt]}, spacecadet, line width=0.9pt]
    (axis cs:-0.52,-0.14) -- (axis cs:-0.52,-0.86);
  \draw[-{Stealth[length=5pt,width=3.5pt]}, spacecadet, line width=0.9pt]
    (axis cs:-0.52,-1.14) -- (axis cs:-0.52,-1.86);
  \node[lbl, anchor=east] at (axis cs:-0.60,-1.5) {$\widehat{L}_-$};

  % Where the ladder breaks off: the next rung would need |m| > sqrt6.
  \draw[-{Stealth[length=5pt,width=3.5pt]}, birrengray, line width=0.9pt,
        densely dashed] (axis cs:0.42,2.14) -- (axis cs:0.42,2.70);
  \draw[birrengray, line width=0.9pt]
    ([xshift=-1.4mm,yshift=-1.4mm]axis cs:0.42,2.70) --
    ([xshift=1.4mm,yshift=1.4mm]axis cs:0.42,2.70);
  \draw[birrengray, line width=0.9pt]
    ([xshift=-1.4mm,yshift=1.4mm]axis cs:0.42,2.70) --
    ([xshift=1.4mm,yshift=-1.4mm]axis cs:0.42,2.70);

  \draw[-{Stealth[length=5pt,width=3.5pt]}, birrengray, line width=0.9pt,
        densely dashed] (axis cs:-0.52,-2.14) -- (axis cs:-0.52,-2.70);
  \draw[birrengray, line width=0.9pt]
    ([xshift=-1.4mm,yshift=-1.4mm]axis cs:-0.52,-2.70) --
    ([xshift=1.4mm,yshift=1.4mm]axis cs:-0.52,-2.70);
  \draw[birrengray, line width=0.9pt]
    ([xshift=-1.4mm,yshift=1.4mm]axis cs:-0.52,-2.70) --
    ([xshift=1.4mm,yshift=-1.4mm]axis cs:-0.52,-2.70);

  \node[lbl, anchor=south] at (axis cs:0.42,3.24)
    {$\widehat{L}_+\ket{\ell,\ell}=0$};
  \node[lbl, anchor=north] at (axis cs:-0.52,-3.24)
    {$\widehat{L}_-\ket{\ell,-\ell}=0$};

  % What the grey rules are
  \node[lbl, font=\scriptsize, anchor=south west, fill=white,
        fill opacity=1, inner sep=2pt]
    at (axis cs:-1.73,2.53) {$|\mathbf{L}|/\hbar=\sqrt{\ell(\ell+1)}=\sqrt{6}$};

  % The multiplet is 2ell+1 states long
  \draw[{Stealth[length=4pt]}-{Stealth[length=4pt]}, birrengray,
        line width=0.6pt] (axis cs:0.95,-2) -- (axis cs:0.95,2);
  \node[lbl, font=\scriptsize, text=birrengray, rotate=90, anchor=south]
    at (axis cs:1.05,0) {$2\ell+1$ \tpcifenglish{states}{Zustände}};
\end{axis}

% ─────────────────────────────────────────────────────────────
%  (b)  The vector model: |L| fixed, L_3 quantised, L_1,L_2 undecided
% ─────────────────────────────────────────────────────────────
\begin{axis}[
    name=cones,
    tpclevelplot,
    scale only axis,
    width=4.78cm, height=5.8cm,      % 6.1 x 7.4 units, equal unit lengths
    at={(ladder.south east)}, anchor=south west, xshift=14mm,
    xmin=-3.05, xmax=3.05,
    ymin=-3.9,  ymax=3.5,
    title={(b) \tpcifenglish{vector model}{Vektormodell}},
]

  % The cones: |L| = sqrt6 hbar with L_3 = m hbar fixed, so the tip runs
  % on a circle of radius sqrt(6 - m^2) at height m.
  \addplot[line width=0.6pt, densely dotted, birrenred,
           domain=0:360, samples=121] ({1.41421*cos(x)}, {2+0.21213*sin(x)});
  \addplot[line width=0.6pt, densely dotted, birrenorange,
           domain=0:360, samples=121] ({2.23607*cos(x)}, {1+0.33541*sin(x)});
  \addplot[line width=0.6pt, densely dotted, munsell,
           domain=0:360, samples=121] ({2.44949*cos(x)}, {0.36742*sin(x)});
  \addplot[line width=0.6pt, densely dotted, cgblue,
           domain=0:360, samples=121] ({2.23607*cos(x)}, {-1+0.33541*sin(x)});
  \addplot[line width=0.6pt, densely dotted, spacecadet,
           domain=0:360, samples=121] ({1.41421*cos(x)}, {-2+0.21213*sin(x)});

  % The sphere of fixed length: every arrow below ends on it.
  \addplot[levellinegrey, domain=0:360, samples=181]
    ({2.44949*cos(x)}, {2.44949*sin(x)});

  % The axis the multiplet is quantised along
  \draw[axisstyle] (axis cs:0,-2.66) -- (axis cs:0,3.30);
  \node[lbl, anchor=west] at (axis cs:0.10,3.22) {$L_3/\hbar$};

  % The second generator of the top cone, so it reads as a cone
  \draw[line width=0.6pt, densely dotted, birrenred]
    (axis cs:0,0) -- (axis cs:-1.41421,2);

  % One representative vector per cone, ending on the circle at
  % (sqrt(6-m^2), m).  Colour = m, as in panel (a).
  \draw[-{Stealth[length=5pt,width=3.5pt]}, birrenred, line width=1.0pt]
    (axis cs:0,0) -- (axis cs:1.41421,2);
  \draw[-{Stealth[length=5pt,width=3.5pt]}, birrenorange, line width=1.0pt]
    (axis cs:0,0) -- (axis cs:2.23607,1);
  \draw[-{Stealth[length=5pt,width=3.5pt]}, munsell, line width=1.0pt]
    (axis cs:0,0) -- (axis cs:2.44949,0);
  \draw[-{Stealth[length=5pt,width=3.5pt]}, cgblue, line width=1.0pt]
    (axis cs:0,0) -- (axis cs:2.23607,-1);
  \draw[-{Stealth[length=5pt,width=3.5pt]}, spacecadet, line width=1.0pt]
    (axis cs:0,0) -- (axis cs:1.41421,-2);

  \node[lbl, anchor=west] at (axis cs:1.52,2)  {$m=\ell$};
  \node[lbl, anchor=west] at (axis cs:1.52,-2) {$m=-\ell$};

  % The tilt that never vanishes: cos theta_min = ell/sqrt(ell(ell+1)),
  % drawn in the gap between the m = 1 and m = 2 cones.
  \addplot[birrengray, line width=0.7pt, domain=54.7356:90, samples=40]
    ({1.68*cos(x)}, {1.68*sin(x)});
  \node[lbl, font=\scriptsize, anchor=west] at (axis cs:0.07,1.13)
    {$\theta_{\min}$};

  % Its pole: the length the top cone falls short of, sqrt6 against m = 2
  \draw[birrengray, line width=0.5pt]
    (axis cs:-0.09,2.4495) -- (axis cs:0.09,2.4495);
  \node[lbl, font=\scriptsize, text=birrengray, anchor=south east,
        inner sep=1pt] at (axis cs:-0.11,2.53) {$\sqrt{6}$};

  % What the dashed circle is
  \node[lbl, anchor=north] at (axis cs:0,-2.72)
    {$|\mathbf{L}|=\hbar\sqrt{\ell(\ell+1)}=\sqrt{6}\,\hbar$};

  % The classical limit as a number, not as an assertion
  \node[lbl, font=\scriptsize, text=birrengray, anchor=north, align=center]
    at (axis cs:0,-3.28)
    {\tpcifenglish{never along the axis}{nie entlang der Achse}:\\[-1pt]
     $\ell=2$: $35^\circ$\quad $\ell=10$: $18^\circ$\quad
     $\ell=100$: $6^\circ$};
\end{axis}

\end{tikzpicture}
\end{center}


Note in (a) the gap between the topmost rung $m=\ell$ and the
dashed bound $\sqrt{\ell(\ell+1)}$: that is precisely where the ladder breaks off.
In (b) the same gap is an angle -- the vector never lies along the axis.

\begin{merksatz}
	Position and momentum are not quantised, angular momentum is: there is a smallest step, of size $\hbar$. Invisible in everyday life, but demonstrated spectroscopically.
\end{merksatz}

We shall now derive the position representation of the eigenfunctions. With the help of the spherical coordinates $r$, $\theta$, $\phi$ we obtain (exercise)
\begin{equation}
	\begin{split}
		\langle r,\theta,\phi|\widehat{L}_3|\Psi\rangle &= -i\hbar \frac{\partial}{\partial \phi} \Psi(r,\theta,\phi), \\
		\langle r,\theta,\phi|\widehat{\mathbf{L}}^2|\Psi\rangle &= -\hbar^2\frac{1}{\sin^2(\theta)}\left(\sin(\theta)\frac{\partial}{\partial \theta}\left(\sin(\theta)\frac{\partial}{\partial \theta}\right) + \frac{\partial^2}{\partial \phi^2}\right) \Psi(r,\theta,\phi),
	\end{split}
\end{equation}
where $|\Psi\rangle \in L^2(\mathbb{R}^3)$ is an arbitrary pure state. 

The eigenfunctions of $\widehat{L}_3$ are given by $e^{im\phi}$:
\begin{equation}
	-i\hbar \frac{\partial}{\partial \phi}e^{im\phi} = \hbar m e^{im\phi}.
\end{equation}
Single-valuedness of the wave function $\Psi(r,\theta,\phi)$ under the shift $\phi\mapsto \phi+2\pi$ demands $m\in \mathbb{Z}$, so that only $\ell = 0,1,2, \ldots$ occurs. (The eigenvectors for half-integer $\ell$ will be constructed later by a different method.) The common eigenfunctions of $\widehat{\mathbf{L}}^2$ and $\widehat{L}_3$ read
\begin{equation}
	\eqbox{\text{Y}_\ell^m (\theta,\phi) \equiv (-1)^m\left(\frac{2\ell + 1}{4\pi}\frac{(\ell-m)!}{(\ell+m)!}\right)^{\frac12} \text{P}_\ell^m(\cos(\theta))e^{im\phi},}
\end{equation}
where, for $m\ge 0$,
\begin{equation}
	\eqbox{\text{P}^m_\ell(x) = (1-x^2)^{\frac{m}{2}}\frac{d^m}{dx^m}\text{P}_\ell(x)}
\end{equation}
are the \emph{associated Legendre functions}, which are defined with the help of the \emph{Legendre polynomials}
\begin{equation}
	\eqbox{\text{P}_\ell(x)=\frac{1}{2^\ell \ell!} \frac{d^\ell}{dx^\ell}(x^2-1)^\ell}
\end{equation}
The functions $\text{Y}_\ell^m$ are known as the \emph{spherical harmonics}. They do not depend on $r$ -- the angular momentum does not see the radial component, and that is what makes the separation ansatz work later on.

The first three spherical harmonics are
\begin{align}
	\text{Y}^0_0(\theta,\phi) &= \frac{1}{\sqrt{4\pi}}, \\
	\text{Y}^0_1(\theta,\phi) &= \sqrt{\frac{3}{4\pi}}\cos(\theta), \\
	\text{Y}^1_1(\theta,\phi) &= -\sqrt{\frac{3}{8\pi}}\sin(\theta)e^{i\phi}.
\end{align}
$\text{Y}^0_0$ is constant: the particle then has no angular momentum at all.
